% GENERATED FILE. Edit canonical lessons, then run npm run generate:books.
% canonical-source-hash: fnv1a64:b8ed6cdd3a62b138
% canonical-lessons: TA-C523-tiruku, TA-C523-mattu, TA-C523-ataippu, TA-C523-vatikal, TA-C523-vatikattu

\chapter{Small Repairs}
\label{ch:ta-small-repairs}
\hlchaptermodality{523}

\begin{chapteropening}
\textbf{By the end of this chapter:} \emph{I can talk about fixing small things at home.}

Complete the last lesson of chapter 523: I can talk about fixing small things at home.
\end{chapteropening}

\section[\texorpdfstring{\ta{திருகு} (tiruku)}{திருகு (tiruku)}]{\ta{திருகு} (tiruku) — to screw; to twist}

\label{lesson:TA-C523-tiruku}



\begin{quote}
Before the new one: say the Tamil for a pillar, then the Tamil for a gutter.
\end{quote}

\subsection*{You'll want to know: \ta{திருகு}}
\textbf{\ta{திருகு}} — \emph{tiruku} — \textquotedblleft{}to screw; to twist\textquotedblright{}.

\begin{quote}
\ta{அவன்} \ta{ஒரு} \ta{புதிய} \ta{குமிழைக்} \ta{கதவில்} \ta{திருகினான்}.
\end{quote}

\begin{quote}
\emph{avaṉ oru putiya kumiḻaik katavil tirukiṉāṉ.}
\end{quote}

\begin{quote}
He screwed a new knob onto the door.
\end{quote}

Past (I …): \ta{திருகினேன்}, \emph{tirukiṉēṉ}.

The same word is a noun: \ta{ஒரு} \ta{திருகு}, a screw.

\begin{cousinweb}
\textbf{\ta{திருகு}} (\emph{tiruku}) is native Dravidian: DEDR 3246, under \ta{திரி} (\emph{tiri}), \textquotedblleft{}to turn\textquotedblright{}, gives \ta{திருகு} (\emph{tiruku}), \textquotedblleft{}to twist, to wring\textquotedblright{}, and the noun \textquotedblleft{}a twist, a screw thread\textquotedblright{}.
\end{cousinweb}

\begin{grammarlens}[title={What you've built}]
One more word for describing what you lived and felt.
\end{grammarlens}

\subsection*{Your turn}
\begin{itemize}
\raggedright
  \item \emph{Say it:} \emph{tiruku}
  \item \emph{Say it:} \emph{tiruku}, once more
  \item \emph{Say it:} \emph{avaṉ oru putiya kumiḻaik katavil tirukiṉāṉ.} — He screwed a new knob onto the door.
  \item \emph{Recall:} say the Tamil for a pillar, then the Tamil for a gutter, then say \emph{tiruku} again
\end{itemize}

\begin{tcolorbox}[breakable,colback=teal!4,colframe=teal!35!black,title={Before you move on}]
What does \ta{திருகு} mean? (\textquotedblleft{}To screw; to twist\textquotedblright{}.) Say it once more.
\end{tcolorbox}
\section[\texorpdfstring{\ta{மாட்டு} (māṭṭu)}{மாட்டு (māṭṭu)}]{\ta{மாட்டு} (māṭṭu) — to hang up; to fix on}

\label{lesson:TA-C523-mattu}



\begin{quote}
Before the new one: say the Tamil for a gutter, then the Tamil for to screw.
\end{quote}

\subsection*{You'll want to know: \ta{மாட்டு}}
\textbf{\ta{மாட்டு}} — \emph{māṭṭu} — \textquotedblleft{}to hang up; to fix on\textquotedblright{}.

\begin{quote}
\ta{அப்பா} \ta{சுவரில்} \ta{ஒரு} \ta{புதிய} \ta{கடிகாரத்தை} \ta{மாட்டினார்}.
\end{quote}

\begin{quote}
\emph{appā cuvaril oru putiya kaṭikārattai māṭṭiṉār.}
\end{quote}

\begin{quote}
Father hung a new clock on the wall.
\end{quote}

Past (I …): \ta{மாட்டினேன்}, \emph{māṭṭiṉēṉ}.

Where it hangs takes -\ta{இல்}: \ta{சுவரில்} \ta{மாட்டு}. Compare \ta{தொங்கு} (\emph{toṅku}), to hang (by itself).

\begin{cousinweb}
\textbf{\ta{மாட்டு}} (\emph{māṭṭu}) is native Dravidian: DEDR 4802 gives \ta{மாட்டு} (\emph{māṭṭu}), \textquotedblleft{}to fasten on, to hook, to fix, to attach\textquotedblright{}.
\end{cousinweb}

\begin{grammarlens}[title={What you've built}]
One more word for describing what you lived and felt.
\end{grammarlens}

\subsection*{Your turn}
\begin{itemize}
\raggedright
  \item \emph{Say it:} \emph{māṭṭu}
  \item \emph{Say it:} \emph{māṭṭu}, once more
  \item \emph{Say it:} \emph{appā cuvaril oru putiya kaṭikārattai māṭṭiṉār.} — Father hung a new clock on the wall.
  \item \emph{Recall:} say the Tamil for a gutter, then the Tamil for to screw, then say \emph{māṭṭu} again
\end{itemize}

\begin{tcolorbox}[breakable,colback=teal!4,colframe=teal!35!black,title={Before you move on}]
What does \ta{மாட்டு} mean? (\textquotedblleft{}To hang up; to fix on\textquotedblright{}.) Say it once more.
\end{tcolorbox}
\section[\texorpdfstring{\ta{அடைப்பு} (aṭaippu)}{அடைப்பு (aṭaippu)}]{\ta{அடைப்பு} (aṭaippu) — a blockage}

\label{lesson:TA-C523-ataippu}



\begin{quote}
Before the new one: say the Tamil for to screw, then the Tamil for to hang up.
\end{quote}

\subsection*{You'll want to know: \ta{அடைப்பு}}
\textbf{\ta{அடைப்பு}} — \emph{aṭaippu} — \textquotedblleft{}a blockage\textquotedblright{}.

\begin{quote}
\ta{குழாயில்} \ta{அடைப்பு} \ta{இருப்பதால்} \ta{தண்ணீர்} \ta{வரவில்லை}.
\end{quote}

\begin{quote}
\emph{kuḻāyil aṭaippu iruppatāl taṇṇīr varavillai.}
\end{quote}

\begin{quote}
Because there is a blockage in the pipe, no water is coming.
\end{quote}

From \ta{அடை}, to shut, to block. A blocked nose is \ta{மூக்கடைப்பு}.

\begin{cousinweb}
\textbf{\ta{அடைப்பு}} (\emph{aṭaippu}) is native Dravidian: DEDR 83 gives \ta{அடை} (\emph{aṭai}), \textquotedblleft{}to shut, to block a passage\textquotedblright{}, and \ta{அடைப்பு}, \textquotedblleft{}closing, an obstruction, a plug\textquotedblright{}.
\end{cousinweb}

\begin{grammarlens}[title={What you've built}]
One more word for describing what you lived and felt.
\end{grammarlens}

\subsection*{Your turn}
\begin{itemize}
\raggedright
  \item \emph{Say it:} \emph{aṭaippu}
  \item \emph{Say it:} \emph{aṭaippu}, once more
  \item \emph{Say it:} \emph{kuḻāyil aṭaippu iruppatāl taṇṇīr varavillai.} — Because there is a blockage in the pipe, no water is coming.
  \item \emph{Recall:} say the Tamil for to screw, then the Tamil for to hang up, then say \emph{aṭaippu} again
\end{itemize}

\begin{tcolorbox}[breakable,colback=teal!4,colframe=teal!35!black,title={Before you move on}]
What does \ta{அடைப்பு} mean? (\textquotedblleft{}A blockage\textquotedblright{}.) Say it once more.
\end{tcolorbox}
\section[\texorpdfstring{\ta{வடிகால்} (vaṭikāl)}{வடிகால் (vaṭikāl)}]{\ta{வடிகால்} (vaṭikāl) — a drain}

\label{lesson:TA-C523-vatikal}



\begin{quote}
Before the new one: say the Tamil for to hang up, then the Tamil for a blockage.
\end{quote}

\subsection*{You'll want to know: \ta{வடிகால்}}
\textbf{\ta{வடிகால்}} — \emph{vaṭikāl} — \textquotedblleft{}a drain\textquotedblright{}.

\begin{quote}
\ta{வடிகால்} \ta{சரியாக} \ta{இல்லாததால்} \ta{வயலில்} \ta{தண்ணீர்} \ta{தேங்கியது}.
\end{quote}

\begin{quote}
\emph{vaṭikāl cariyāka illātatāl vayalil taṇṇīr tēṅkiyatu.}
\end{quote}

\begin{quote}
Because the drain was not right, water stood in the field.
\end{quote}

\ta{வடி} \textquotedblleft{}to drain away\textquotedblright{} + \ta{கால்} \textquotedblleft{}a channel\textquotedblright{}.

\begin{cousinweb}
\textbf{\ta{வடிகால்}} (\emph{vaṭikāl}) is a native compound: \ta{வடி} (\emph{vaṭi}), \textquotedblleft{}to drip, to flow out\textquotedblright{} (DEDR 5221), and \ta{கால்} (\emph{kāl}), \textquotedblleft{}a channel\textquotedblright{}.
\end{cousinweb}

\begin{grammarlens}[title={What you've built}]
One more word for describing what you lived and felt.
\end{grammarlens}

\subsection*{Your turn}
\begin{itemize}
\raggedright
  \item \emph{Say it:} \emph{vaṭikāl}
  \item \emph{Say it:} \emph{vaṭikāl}, once more
  \item \emph{Say it:} \emph{vaṭikāl cariyāka illātatāl vayalil taṇṇīr tēṅkiyatu.} — Because the drain was not right, water stood in the field.
  \item \emph{Recall:} say the Tamil for to hang up, then the Tamil for a blockage, then say \emph{vaṭikāl} again
\end{itemize}

\begin{tcolorbox}[breakable,colback=teal!4,colframe=teal!35!black,title={Before you move on}]
What does \ta{வடிகால்} mean? (\textquotedblleft{}A drain\textquotedblright{}.) Say it once more.
\end{tcolorbox}
\section[\texorpdfstring{\ta{வடிகட்டு} (vaṭikaṭṭu)}{வடிகட்டு (vaṭikaṭṭu)}]{\ta{வடிகட்டு} (vaṭikaṭṭu) — to filter; to strain}

\label{lesson:TA-C523-vatikattu}



\begin{quote}
Before the new one: say the Tamil for a blockage, then the Tamil for a drain.
\end{quote}

\subsection*{You'll want to know: \ta{வடிகட்டு}}
\textbf{\ta{வடிகட்டு}} — \emph{vaṭikaṭṭu} — \textquotedblleft{}to filter; to strain\textquotedblright{}.

\begin{quote}
\ta{அம்மா} \ta{தேநீரை} \ta{ஒரு} \ta{துணியால்} \ta{வடிகட்டினாள்}.
\end{quote}

\begin{quote}
\emph{ammā tēnīrai oru tuṇiyāl vaṭikaṭṭiṉāḷ.}
\end{quote}

\begin{quote}
Mother strained the tea through a cloth.
\end{quote}

Past (I …): \ta{வடிகட்டினேன்}, \emph{vaṭikaṭṭiṉēṉ}.

What is strained takes -\ta{ஐ}; the cloth or sieve takes -\ta{ஆல்}, \textquotedblleft{}with\textquotedblright{}.

\begin{cousinweb}
\textbf{\ta{வடிகட்டு}} (\emph{vaṭikaṭṭu}) is a native compound built on \ta{வடி} (\emph{vaṭi}), \textquotedblleft{}to drip, to cause to flow out, to strain\textquotedblright{} (DEDR 5221).
\end{cousinweb}

\begin{grammarlens}[title={What you've built}]
One more word for describing what you lived and felt.
\end{grammarlens}

\subsection*{Your turn}
\begin{itemize}
\raggedright
  \item \emph{Say it:} \emph{vaṭikaṭṭu}
  \item \emph{Say it:} \emph{vaṭikaṭṭu}, once more
  \item \emph{Say it:} \emph{ammā tēnīrai oru tuṇiyāl vaṭikaṭṭiṉāḷ.} — Mother strained the tea through a cloth.
  \item \emph{Recall:} say the Tamil for a blockage, then the Tamil for a drain, then say \emph{vaṭikaṭṭu} again
\end{itemize}

\begin{tcolorbox}[breakable,colback=teal!4,colframe=teal!35!black,title={Before you move on}]
What does \ta{வடிகட்டு} mean? (\textquotedblleft{}To filter; to strain\textquotedblright{}.) Say it once more.
\end{tcolorbox}
